\documentclass[11pt,a4paper]{article}

% --- Packages ---
\usepackage[utf8]{inputenc}
\usepackage[margin=1in]{geometry}
\usepackage{amsmath,amssymb,amsfonts}
\usepackage{booktabs}
\usepackage{microtype}
\usepackage{xcolor}
\usepackage{hyperref}

\hypersetup{
    colorlinks=true,
    linkcolor=blue!70!black,
    citecolor=blue!70!black,
    urlcolor=blue!70!black
}

\title{\textbf{MATH502: Numerical Methods}\\[0.3em]
\Large Detailed Hand Analysis: Problem 1(a) \& Problem 6(a)}
\author{\textbf{Mazen Altelbany} (ID: 64-12371, Tutorial: T-04) \\
\textbf{Yahia Talaat} (ID: 64-3581, Tutorial: T-02)}
\date{Winter 2026 Semester}

\begin{document}

\maketitle

\hrule
\vspace{1em}

% =========================================================================
% SECTION 1: PROBLEM 1(a) HAND ANALYSIS
% =========================================================================
\section{Problem 1(a): Pen-and-Paper Manual Iterations for LORAN 2D Navigation}

\subsection{Non-linear System Definition}
The LOng RAnge Navigation (LORAN) receiver determines the vessel's coordinates $(x, y)$ by finding the intersection of two hyperbolic lines of position defined by time difference of arrival (TDOA) signals:
\begin{align}
    f_1(x, y) &= \frac{x^2}{186^2} - \frac{y^2}{300^2 - 186^2} - 1 = \frac{x^2}{34596} - \frac{y^2}{55404} - 1 = 0, \label{eq:f1} \\
    f_2(x, y) &= \frac{(y - 500)^2}{279^2} - \frac{(x - 300)^2}{500^2 - 279^2} - 1 = \frac{(y - 500)^2}{77841} - \frac{(x - 300)^2}{172159} - 1 = 0. \label{eq:f2}
\end{align}

\subsection{Analytical Jacobian Matrix}
The multi-dimensional Newton-Raphson method requires the Jacobian matrix $\mathbf{J}(x, y) \in \mathbb{R}^{2 \times 2}$:
\begin{equation}
    \mathbf{J}(x, y) = 
    \begin{bmatrix}
        \dfrac{\partial f_1}{\partial x} & \dfrac{\partial f_1}{\partial y} \\[1em]
        \dfrac{\partial f_2}{\partial x} & \dfrac{\partial f_2}{\partial y}
    \end{bmatrix}
    =
    \begin{bmatrix}
        \dfrac{2x}{34596} & -\dfrac{2y}{55404} \\[1em]
        -\dfrac{2(x - 300)}{172159} & \dfrac{2(y - 500)}{77841}
    \end{bmatrix}.
\end{equation}

At each iteration $k$, the Newton correction vector $\Delta \mathbf{x}_k = [\Delta x_k, \Delta y_k]^T$ is obtained by solving:
\begin{equation}
    \mathbf{J}(\mathbf{x}_k) \Delta \mathbf{x}_k = -F(\mathbf{x}_k) = 
    \begin{bmatrix}
        -f_1(x_k, y_k) \\
        -f_2(x_k, y_k)
    \end{bmatrix},
\end{equation}
followed by the coordinate update:
\begin{equation}
    \mathbf{x}_{k+1} = \mathbf{x}_k + \Delta \mathbf{x}_k = 
    \begin{bmatrix}
        x_k + \Delta x_k \\
        y_k + \Delta y_k
    \end{bmatrix}.
\end{equation}

% -------------------------------------------------------------------------
\subsection{Iteration 1: Step-by-Step Hand Calculations}
Starting from the given initial estimate $\mathbf{x}_0 = [x_0, y_0]^T = [200, 200]^T$:

\paragraph{Step 1: Evaluate the Residual Vector $F(\mathbf{x}_0)$:}
\begin{align*}
    f_1(200, 200) &= \frac{200^2}{34596} - \frac{200^2}{55404} - 1 = 1.15620303 - 0.72196953 - 1 = -0.56576650, \\
    f_2(200, 200) &= \frac{(200 - 500)^2}{77841} - \frac{(200 - 300)^2}{172159} - 1 = \frac{90000}{77841} - \frac{10000}{172159} - 1 \\
                  &= 1.15620303 - 0.05808584 - 1 = +0.09811719.
\end{align*}
Thus, the right-hand side vector $-F(\mathbf{x}_0)$ is:
\begin{equation*}
    -F(\mathbf{x}_0) = 
    \begin{bmatrix}
        -f_1(200, 200) \\
        -f_2(200, 200)
    \end{bmatrix}
    =
    \begin{bmatrix}
        +0.56576650 \\
        -0.09811719
    \end{bmatrix}.
\end{equation*}

\paragraph{Step 2: Evaluate the Jacobian Matrix $\mathbf{J}_0$:}
\begin{align*}
    J_{11} &= \frac{2(200)}{34596} = \frac{400}{34596} \approx 0.01156203, \\
    J_{12} &= -\frac{2(200)}{55404} = -\frac{400}{55404} \approx -0.00721970, \\
    J_{21} &= -\frac{2(200 - 300)}{172159} = \frac{200}{172159} \approx 0.00116172, \\
    J_{22} &= \frac{2(200 - 500)}{77841} = -\frac{600}{77841} \approx -0.00770802.
\end{align*}
\begin{equation*}
    \mathbf{J}_0 = 
    \begin{bmatrix}
        0.01156203 & -0.00721970 \\
        0.00116172 & -0.00770802
    \end{bmatrix}.
\end{equation*}

\paragraph{Step 3: Compute the Jacobian Determinant $\det(\mathbf{J}_0)$:}
\begin{align*}
    \det(\mathbf{J}_0) &= J_{11} J_{22} - J_{12} J_{21} \\
    &= (0.01156203)(-0.00770802) - (-0.00721970)(0.00116172) \\
    &= -8.912039 \times 10^{-5} - (-8.387271 \times 10^{-6}) \\
    &= -8.073312 \times 10^{-5}.
\end{align*}

\paragraph{Step 4: Solve for Correction $\Delta \mathbf{x}_0$ via Cramer's Rule:}
\begin{align*}
    \Delta x_0 &= \frac{\det \begin{bmatrix} -f_1 & J_{12} \\ -f_2 & J_{22} \end{bmatrix}}{\det(\mathbf{J}_0)} 
    = \frac{(+0.56576650)(-0.00770802) - (-0.09811719)(-0.00721970)}{-8.073312 \times 10^{-5}} \\
    &= \frac{-0.00436094 - 0.00070838}{-8.073312 \times 10^{-5}} = \frac{-0.00506932}{-8.073312 \times 10^{-5}} = +62.791030. \\[1em]
    \Delta y_0 &= \frac{\det \begin{bmatrix} J_{11} & -f_1 \\ J_{21} & -f_2 \end{bmatrix}}{\det(\mathbf{J}_0)} 
    = \frac{(0.01156203)(-0.09811719) - (0.00116172)(+0.56576650)}{-8.073312 \times 10^{-5}} \\
    &= \frac{-0.00113444 - 0.00065726}{-8.073312 \times 10^{-5}} = \frac{-0.00179170}{-8.073312 \times 10^{-5}} = +22.192804.
\end{align*}

\paragraph{Step 5: Coordinate Update after Iteration 1:}
\begin{align*}
    x_1 &= x_0 + \Delta x_0 = 200.0 + 62.791030 = 262.791030, \\
    y_1 &= y_0 + \Delta y_0 = 200.0 + 22.192804 = 222.192804.
\end{align*}

% -------------------------------------------------------------------------
\subsection{Iteration 2: Step-by-Step Hand Calculations}
Starting from $\mathbf{x}_1 = [262.791030, 222.192804]^T$:

\paragraph{Step 1: Evaluate the Residual Vector $F(\mathbf{x}_1)$:}
\begin{align*}
    f_1(262.791030, 222.192804) &= \frac{262.791030^2}{34596} - \frac{222.192804^2}{55404} - 1 \\
    &= 1.99616086 - 0.89108544 - 1 = +0.10507542, \\
    f_2(262.791030, 222.192804) &= \frac{(222.192804 - 500)^2}{77841} - \frac{(262.791030 - 300)^2}{172159} - 1 \\
    &= \frac{(-277.807196)^2}{77841} - \frac{(-37.208970)^2}{172159} - 1 \\
    &= 0.99146914 - 0.00804204 - 1 = -0.01657290.
\end{align*}
Hence:
\begin{equation*}
    -F(\mathbf{x}_1) = 
    \begin{bmatrix}
        -0.10507542 \\
        +0.01657290
    \end{bmatrix}.
\end{equation*}

\paragraph{Step 2: Evaluate the Jacobian Matrix $\mathbf{J}_1$:}
\begin{align*}
    J_{11} &= \frac{2(262.791030)}{34596} \approx 0.01519199, \\
    J_{12} &= -\frac{2(222.192804)}{55404} \approx -0.00802082, \\
    J_{21} &= -\frac{2(262.791030 - 300)}{172159} = \frac{2(37.208970)}{172159} \approx 0.00043226, \\
    J_{22} &= \frac{2(222.192804 - 500)}{77841} = -\frac{2(277.807196)}{77841} \approx -0.00713781.
\end{align*}
\begin{equation*}
    \mathbf{J}_1 = 
    \begin{bmatrix}
        0.01519199 & -0.00802082 \\
        0.00043226 & -0.00713781
    \end{bmatrix}.
\end{equation*}

\paragraph{Step 3: Compute the Jacobian Determinant $\det(\mathbf{J}_1)$:}
\begin{align*}
    \det(\mathbf{J}_1) &= J_{11} J_{22} - J_{12} J_{21} \\
    &= (0.01519199)(-0.00713781) - (-0.00802082)(0.00043226) \\
    &= -1.084375 \times 10^{-4} - (-3.467080 \times 10^{-6}) \\
    &= -1.049704 \times 10^{-4}.
\end{align*}

\paragraph{Step 4: Solve for Correction $\Delta \mathbf{x}_1$ via Cramer's Rule:}
\begin{align*}
    \Delta x_1 &= \frac{\det \begin{bmatrix} -f_1 & J_{12} \\ -f_2 & J_{22} \end{bmatrix}}{\det(\mathbf{J}_1)} 
    = \frac{(-0.10507542)(-0.00713781) - (+0.01657290)(-0.00802082)}{-1.049704 \times 10^{-4}} \\
    &= \frac{0.00075001 - (-0.00013293)}{-1.049704 \times 10^{-4}} = \frac{0.00088294}{-1.049704 \times 10^{-4}} = -8.411356. \\[1em]
    \Delta y_1 &= \frac{\det \begin{bmatrix} J_{11} & -f_1 \\ J_{21} & -f_2 \end{bmatrix}}{\det(\mathbf{J}_1)} 
    = \frac{(0.01519199)(+0.01657290) - (0.00043226)(-0.10507542)}{-1.049704 \times 10^{-4}} \\
    &= \frac{0.00025178 - (-0.00004542)}{-1.049704 \times 10^{-4}} = \frac{0.00029720}{-1.049704 \times 10^{-4}} = -2.831433.
\end{align*}

\paragraph{Step 5: Coordinate Update after Iteration 2:}
\begin{align*}
    x_2 &= x_1 + \Delta x_1 = 262.791030 - 8.411356 = 254.379674, \\
    y_2 &= y_1 + \Delta y_1 = 222.192804 - 2.831433 = 219.361371.
\end{align*}

\subsection{Summary of Manual Iterations vs. Converged Solution}
\begin{table}[h]
\centering
\caption{Manual Iterations compared to Final Converged Root ($tol = 10^{-8}$).}
\label{tab:p1_summary}
\begin{tabular}{cccccc}
\toprule
\textbf{Iteration $k$} & $x_k$ & $y_k$ & $\|F(\mathbf{x}_k)\|_2$ & $\|\Delta \mathbf{x}_k\|_2$ & Status \\
\midrule
$0$ (Initial) & $200.000000$ & $200.000000$ & $5.742 \times 10^{-1}$ & $6.660 \times 10^{+1}$ & Manual step \\
$1$           & $262.791030$ & $222.192804$ & $1.064 \times 10^{-1}$ & $8.875 \times 10^{0}$  & Manual step \\
$2$           & $254.379674$ & $219.361371$ & $2.201 \times 10^{-3}$ & $1.670 \times 10^{-1}$ & Manual step \\
\midrule
$4$ (Converged) & \textbf{254.221120} & \textbf{219.306992} & \textbf{$8.546 \times 10^{-14}$} & \textbf{$7.430 \times 10^{-12}$} & Converged \\
\bottomrule
\end{tabular}
\end{table}

\vspace{1.5em}
\hrule
\vspace{1.5em}

% =========================================================================
% SECTION 2: PROBLEM 6(a) HAND ANALYSIS
% =========================================================================
\section{Problem 6(a): Analytical Closed-Form Solution Derivation (Weierstrass Substitution)}

\subsection{Physical Context \& Integral Definition}
Consider a ground Wirelessly Chargeable Device (WCD) harvesting RF energy from a constellation of $N$ stacked Low Earth Orbit (LEO) satellites, each equipped with $M$ antennas. 
Applying Maximum Ratio Transmission (MRT) beamforming, the received power at the WCD is:
\begin{equation}
    P_r(t) = M N^2 P_{r,\text{single}}(t) = M N^2 P_t G_t G_r \left( \frac{\lambda}{4\pi d(t)} \right)^2 \alpha_s \beta_s,
\end{equation}
where $d(t)$ is the distance between the constellation grid and the ground WCD:
\begin{equation}
    d(t) = \sqrt{(R + H)^2 + R^2 - 2(R + H)R \cos\phi \cos(\omega t)}.
\end{equation}

The harvested power is $P_h(t) = \eta_h P_r(t) = \dfrac{M N^2 \mu}{d^2(t)}$, where the Friis scaling constant is:
\begin{equation}
    \mu = \eta_h P_t G_t G_r \left( \frac{\lambda}{4\pi} \right)^2 \alpha_s \beta_s.
\end{equation}

As clarified by the course instructor, the factor $M N^2$ is applied directly to the received power and is \textbf{multiplied only once} (no double multiplication). Furthermore, accounting for the full symmetric visibility pass as the satellite constellation approaches from the horizon ($-T/2$), passes zenith ($t=0$), and recedes ($+T/2$), the total harvested energy $E_h$ is:
\begin{equation}
    E_h = \int_{-T/2}^{T/2} P_h(t) \, dt = 2 \int_0^{T_{\text{half}}} P_h(t) \, dt = 2 M N^2 \mu \int_0^{T_{\text{half}}} \frac{dt}{(R+H)^2 + R^2 - 2(R+H)R \cos\phi \cos(\omega t)},
\end{equation}
where $T_{\text{half}} = \theta_c / \omega$.

\subsection{Step-by-Step Derivation via Weierstrass Substitution}

\paragraph{Step 1: Define Constants $A$ and $B$:}
Let:
\begin{align}
    A &= (R + H)^2 + R^2 = H^2 + 2HR + 2R^2, \\
    B &= 2R(R + H)\cos\phi.
\end{align}
Then the integral is written compactly as:
\begin{equation}
    E_h = 2 M N^2 \mu \int_0^{T_{\text{half}}} \frac{dt}{A - B\cos(\omega t)}.
\end{equation}

\paragraph{Step 2: Apply the Weierstrass Half-Angle Substitution ($u$-substitution):}
Let:
\begin{equation}
    u = \tan\left( \frac{\omega t}{2} \right).
\end{equation}
From standard trigonometric half-angle identities:
\begin{equation}
    \cos(\omega t) = \frac{1 - \tan^2(\omega t / 2)}{1 + \tan^2(\omega t / 2)} = \frac{1 - u^2}{1 + u^2}.
\end{equation}
Differentiating $u$ with respect to $t$:
\begin{equation}
    du = \frac{\omega}{2} \sec^2\left(\frac{\omega t}{2}\right) dt = \frac{\omega}{2}(1 + u^2) dt \implies dt = \frac{2}{\omega} \frac{du}{1 + u^2}.
\end{equation}

\paragraph{Step 3: Transform the Integration Limits:}
\begin{itemize}
    \item Lower limit ($t = 0$):
    \begin{equation*}
        u(0) = \tan(0) = 0.
    \end{equation*}
    \item Upper limit ($t = T_{\text{half}}$):
    \begin{equation*}
        u(T_{\text{half}}) = \tan\left( \frac{\omega T_{\text{half}}}{2} \right).
    \end{equation*}
\end{itemize}

\paragraph{Step 4: Substitute into the Integral:}
Substitute $\cos(\omega t)$ and $dt$ into the integrand:
\begin{align*}
    A - B\cos(\omega t) &= A - B\left( \frac{1 - u^2}{1 + u^2} \right) = \frac{A(1 + u^2) - B(1 - u^2)}{1 + u^2} \\
    &= \frac{(A - B) + (A + B)u^2}{1 + u^2}.
\end{align*}
Therefore:
\begin{align}
    E_h &= 2 M N^2 \mu \int_0^{\tan(\omega T_{\text{half}} / 2)} \frac{1}{\dfrac{(A - B) + (A + B)u^2}{1 + u^2}} \cdot \left( \frac{2}{\omega} \frac{du}{1 + u^2} \right) \nonumber \\
        &= \frac{4 M N^2 \mu}{\omega} \int_0^{\tan(\omega T_{\text{half}} / 2)} \frac{du}{(A - B) + (A + B)u^2}.
\end{align}

\paragraph{Step 5: Analytical Evaluation of the Standard Arctangent Form:}
Recall the standard calculus integral for $\alpha^2 > 0$ and $\beta^2 > 0$:
\begin{equation}
    \int \frac{du}{\alpha^2 + \beta^2 u^2} = \frac{1}{\alpha \beta} \tan^{-1}\left( \frac{\beta}{\alpha} u \right).
\end{equation}
Here, $\alpha^2 = A - B$ and $\beta^2 = A + B$. Because orbital altitude $H > 0$ and radius $R > 0$, we have $A > B > 0$. Thus:
\begin{equation}
    \alpha \beta = \sqrt{(A - B)(A + B)} = \sqrt{A^2 - B^2}, \quad \text{and} \quad \frac{\beta}{\alpha} = \sqrt{\frac{A + B}{A - B}}.
\end{equation}
Evaluating between $u = 0$ and $u = \tan(\omega T_{\text{half}} / 2)$:
\begin{equation}
    \int_0^{\tan(\omega T_{\text{half}} / 2)} \frac{du}{(A - B) + (A + B)u^2} = \frac{1}{\sqrt{A^2 - B^2}} \tan^{-1}\left( \sqrt{\frac{A + B}{A - B}} \tan\left(\frac{\omega T_{\text{half}}}{2}\right) \right).
\end{equation}

\paragraph{Step 6: Evaluation for Aligned Ground Track ($\phi = 0$):}
For an aligned satellite track with azimuth angle $\phi = 0$, $\cos\phi = 1$:
\begin{align}
    A - B &= (R + H)^2 + R^2 - 2R(R + H) = \big((R + H) - R\big)^2 = H^2, \\
    A + B &= (R + H)^2 + R^2 + 2R(R + H) = \big((R + H) + R\big)^2 = (H + 2R)^2.
\end{align}
Taking square roots (since $H > 0$ and $R > 0$):
\begin{align}
    \sqrt{A^2 - B^2} &= \sqrt{(A - B)(A + B)} = \sqrt{H^2 (H + 2R)^2} = H(H + 2R), \\
    \sqrt{\frac{A + B}{A - B}} &= \sqrt{\frac{(H + 2R)^2}{H^2}} = \frac{H + 2R}{H}.
\end{align}

\subsection{Final Closed-Form Analytical Expression}
Substituting these geometric expressions back into the energy equation:
\begin{equation}
    \boxed{E_h = \frac{4 M N^2 \mu}{\omega H (H + 2R)} \tan^{-1}\left( \frac{H + 2R}{H} \tan\left( \frac{\omega T_{\text{half}}}{2} \right) \right)}
\end{equation}

\subsection{Equivalence to Dr. Shehab's Cotangent Formulation}
Using the reciprocal trigonometric identity for the arctangent function:
\begin{equation}
    \tan^{-1}(z) = \frac{\pi}{2} - \tan^{-1}\left( \frac{1}{z} \right), \quad \text{for } z > 0,
\end{equation}
substituting $z = \left( \dfrac{H + 2R}{H} \right) \tan\left(\dfrac{\omega T_{\text{half}}}{2}\right)$:
\begin{align*}
    \tan^{-1}\left( \frac{H + 2R}{H} \tan\left(\frac{\omega T_{\text{half}}}{2}\right) \right) 
    &= \frac{\pi}{2} - \tan^{-1}\left( \frac{H}{(H + 2R) \tan(\omega T_{\text{half}} / 2)} \right) \\
    &= \frac{\pi}{2} - \tan^{-1}\left( \frac{H \cot(\omega T_{\text{half}} / 2)}{H + 2R} \right) \\
    &= \frac{1}{2} \left[ \pi - 2 \tan^{-1}\left( \frac{H \cot(\omega T_{\text{half}} / 2)}{H + 2R} \right) \right].
\end{align*}
Multiplying by $\dfrac{4 M N^2 \mu}{\omega H (H + 2R)}$:
\begin{equation}
    \boxed{E_h = \frac{2 M N^2 \mu}{H^2 \omega + 2HR\omega} \left[ \pi - 2 \tan^{-1}\left( \frac{H \cot(\omega T_{\text{half}} / 2)}{H + 2R} \right) \right]}
\end{equation}
This mathematically proves exact equivalence to Dr.~Shehab's published research paper formula (IEEE, 2024), fully accounting for the symmetric visibility region ($2 \times \int_0^{T_{\text{half}}}$) and single $M N^2$ beamforming factor.

\end{document}

